\SourceLeaf{48}
sed quoque antecedentis $(r-1)^{\text{ti}}$ posteriorem
\[
= \frac{1}{2^{2m-1}}\,{}^{2m-1}\mathfrak{B}^{(m+r-1)}\cos 2r\varphi,
\]
esse in $\cos 2r\varphi$ ductam, ex ceteris terminis vero multipla aut minora
aut maiora proficisci. Iunctis igitur terminis hisce habebitur novae seriei
pro $\cos\varphi^{2m}$ terminus generalis $r^{\text{tus}}$, qui, cum sit
\[
{}^{2m-1}\mathfrak{B}^{(m+r)} + {}^{2m-1}\mathfrak{B}^{(m+r-1)} = {}^{2m}\mathfrak{B}^{(m+r)},
\]
induet formam
\[
\frac{1}{2^{2m-1}}\,{}^{2m}\mathfrak{B}^{(m+r)} . \cos 2r\varphi.
\]

\par\noindent Haec expressio legi infra pro potestatibus paribus adsumtae prorsus
consentanea est, ex quo sequitur, valituram illam pro potestati pari, simulac
impari proxime minori demonstrata fuerit.

4. Denuo igitur derivata modo pro $\cos\varphi^{2m}$ series:
\[
\cos\varphi^{2m} = \frac{1}{2^{2m-1}}\frac{{}^{2m}\mathfrak{B}^{(m)}}{2}
+ \frac{1}{2^{2m-1}}\,{}^{2m}\mathfrak{B}^{(m+1)}\cos 2\varphi
+ \frac{1}{2^{2m-1}}\,{}^{2m}\mathfrak{B}^{(m+2)}\cos 4\varphi + .\,.\,.
\]
\[
{} + \frac{1}{2^{2m-1}}\,{}^{2m}\mathfrak{B}^{(m+r)}\cos 2r\varphi + .\,.\,.
\]
in $\cos\varphi$ ducatur, ut habeatur ex una parte $\cos\varphi^{2m+1}$, ex
altera, excepto initiali, in quolibet termino binorum cosinuum productum, ut
antea in binas partes dirimendum. Dabit sic terminus primus
\[
\frac{1}{2^{2m-1}}\,{}^{2m}\mathfrak{B}^{(m+1)}\cos 2\varphi\cos\varphi
\]
partes
\[
\frac{1}{2^{2m-1}}\,{}^{2m}\mathfrak{B}^{(m+1)}\frac{1}{2}\cos\varphi
+ \frac{1}{2^{2m-1}}\,{}^{2m}\mathfrak{B}^{(m+1)}\frac{1}{2}\cos 3\varphi.
\]

\par\noindent Pars eius prima, cum unica, quam terminus initialis praebet
\[
\frac{1}{2^{2m-1}}\frac{{}^{2m}\mathfrak{B}^{(m)}}{2}\cos\varphi,
\]
in summam collecta, dabit novae seriei initialem
\[
\frac{1}{2^{2m}}\,{}^{2m+1}\mathfrak{B}^{(m+1)}\cos\varphi.
\]
